\chapter{Methodology}
\label{ch:methodology}

% Allow a little extra line stretch so long technical compounds rebreak
% cleanly instead of spilling into the margin (no template packages added).
\setlength{\emergencystretch}{3em}

This chapter describes how the project was built rather than what it does. It
presents the development model adopted --- a design-led, AI-assisted single-developer
process (Section~\ref{sec:ai-model}); the sequence of earlier attempts and the
lessons that led to the final approach (Section~\ref{sec:evolution}); the software
process itself, a continuous-flow adaptation of Agile to a solo developer working
with an AI assistant (Section~\ref{sec:process}); the technologies and tooling used
(Section~\ref{sec:tech}); and, as the ESII regulation requires, an explicit
declaration of the use of artificial-intelligence tools
(Section~\ref{sec:ai-declaration}).

\section{A design-led, AI-assisted development model}
\label{sec:ai-model}

The defining methodological choice of this project is the division of labour between
the author and an AI coding assistant. The author acted as \emph{architect and
designer} --- eliciting requirements, deciding the architecture, and reviewing and
validating every output --- while the assistant carried out most of the
implementation under that direction. This inverts the traditional emphasis of a
software project: the scarce effort was spent on design, software-engineering
judgement and architecture, not on typing code.

The evidence for this model is the decision trail. Every significant choice was
recorded as an Architecture Decision Record (ADR) in the lightweight MADR format
\cite{madr, nygard_adr}; the repository contains \textbf{thirty-five} such records,
each stating a context, the options considered, and the decision with its
consequences. Requirements were captured in a Software Requirements Specification
following the structure of ISO/IEC/IEEE~29148 \cite{iso29148}. Because these
artefacts are authored by the developer and precede the code, they demonstrate that
the intellectual work --- what to build and why --- remained a human responsibility,
with the assistant supplying the implementation. It also reflects a disciplined
approach to AI-assisted development: reviewing every change before accepting it,
decomposing work into small increments, letting the assistant run the tests under
supervision, and governing it with a persistent rules file.

\section{Evolution of the approach}
\label{sec:evolution}

The final method was not the first one tried. It emerged from a series of earlier
projects whose shortcomings shaped it, and the progression is itself instructive.

\subsection{From a pentesting platform to a retesting tool}
\label{subsec:prior-attempts}

The project began under the name \emph{Adeptus} with an ambition that proved too
broad: a full, locally deployable \emph{penetration-testing platform}. Successive
iterations explored this idea --- an AI-powered pentesting web application, and then
a multi-user ``pentest workbench'' pairing each operator with a private AI
conversation over a shared, engagement-scoped knowledge graph, extensible through
external tools. The scope was the problem. Building a general offensive platform
meant most of the effort went into breadth --- orchestration, multi-user concerns,
tool integration --- rather than into a defensible, evaluable research
contribution.

The turning point was narrowing the problem from \emph{pentesting} to
\emph{retesting}. A focused prototype, \emph{Retester}, discarded the platform
ambition and did one thing: analyse a PDF report, extract each finding, and let an
AI agent re-verify it against the audited system in an interactive, steerable
session. That prototype is the conceptual seed of the system presented in this
dissertation. Table~\ref{tab:evolution} summarises the progression and the role each
attempt played.

\begin{table}[htb]
\footnotesize
\renewcommand{\arraystretch}{1.3}
\begin{center}
\begin{tabular}{|L{2.6cm}|L{3.4cm}|L{3.0cm}|L{4.0cm}|}
\hline
\rowcolor{tema!10}
\bft{Attempt} & \bft{Scope / focus} & \bft{AI-workflow investment} & \bft{Role in the path to this work} \\ \hline
\emph{Adeptus} (hand-crafted) & AI-powered pentesting platform & Ad-hoc prompting; no skills, hooks or ADRs & First attempt; showed that unstructured AI use does not scale. \\ \hline
\emph{Adeptus} (workbench) & Multi-user pentest workbench + knowledge graph & Heavy: 13 skills, 7 role-agents, 7 hooks, slice-based flow & Peak tooling; revealed that ceremony had outgrown the task. \\ \hline
\emph{Retester} & Report-driven vulnerability \emph{revalidation} & Focused; Pydantic AI, swappable Claude/Ollama & The pivot from pentesting to retesting. \\ \hline
\bft{\emph{revalid}} (this work) & Report-driven, human-gated retest tool & Right-sized: 9 skills, 3 agents, 3 hooks, 35 ADRs, Kanban + CI & The disciplined, evaluable TFG. \\ \hline
\end{tabular}
\end{center}
\caption[Evolution of the project across earlier attempts]{Evolution from the
\emph{Adeptus} pentesting platform to the focused \emph{revalid} retesting tool. The
decisive lessons were to narrow the scope from pentesting to revalidation, and to
right-size the AI-workflow ceremony to the value of a single-developer thesis.}
\label{tab:evolution}
\end{table}

\subsection{From a hybrid process to disciplined iteration}
\label{subsec:process-evolution}

The way of working evolved in parallel with the scope. The earliest attempts mixed
an up-front, waterfall-style plan with ad-hoc agile execution, and the combination
was unwieldy: a large plan invited large, hard-to-review changes, which is precisely
where AI-assisted development goes wrong. Through trial and error the process
settled on strict iteration --- build a minimal but working end-to-end slice first,
then expand it --- which is the walking-skeleton principle applied throughout this
project.

A second lesson concerned tooling. The workbench iteration of \emph{Adeptus} had
accumulated thirteen custom assistant skills, seven role-playing agents and a
seven-hook automation layer; this was more process than a solo thesis could justify.
The present project deliberately \emph{reduced} that apparatus to nine
task-specific skills, three quality-oriented agents and three hooks, and made the
reduction an explicit, recorded decision (ADR-0004, ``ceremony scales with thesis
value''), trading machinery for discipline.

\section{Software process: Kanban for solo, AI-paced development}
\label{sec:process}

AI-assisted development compresses the cycle time of a feature from weeks to hours,
which makes the fixed-length sprints of Scrum \cite{scrumguide} a poor fit. The
project therefore adopted a continuous-flow \textbf{Kanban} process
\cite{kanban_anderson} on a GitHub Projects board, as an explicit adaptation of the
values of the Agile Manifesto \cite{agilemanifesto} to a single developer working
with an assistant. Work flows continuously from a prioritised backlog; there are no
sprints and no meetings; and \emph{milestones} mark feature-complete increments
rather than time boxes. The project passed through six such milestones, from the M1
walking skeleton to the M6 agentic console; the first four each closed with a
semantically versioned release \cite{semver} (v0.1.0 through v0.4.0), while M5 and
M6 were in progress at the time of writing. Figure~\ref{fig:milestones}
lays out this milestone sequence.

\begin{figure}[htb]
\centering
\resizebox{\textwidth}{!}{%
\begin{tikzpicture}[
  >={Stealth[length=2mm]},
  ms/.style={draw=gris50, thick, rounded corners=2pt, align=center, text width=2.3cm, minimum height=1.0cm, fill=gris95, font=\scriptsize, inner sep=3pt},
  rel/.style={font=\scriptsize\bfseries, text=gris25},
]
\draw[->, thick, gris50] (-0.4,0) -- (14.8,0);
\foreach \x/\m/\l in {1/M1/{walking skeleton}, 3.3/M2/{report understanding}, 5.6/M3/{React SPA}, 7.9/M4/{verdicts, audit, export}, 10.2/M5/{evaluation \& config}, 12.5/M6/{agentic console}} {
  \draw[gris50, thick] (\x,0.1) -- (\x,-0.1);
  \node[ms, anchor=south] at (\x,0.35) {\bft{\m}\\\l};
}
\foreach \x/\v in {1/{v0.1.0}, 3.3/{v0.2.0}, 5.6/{v0.3.0}, 7.9/{v0.4.0}} {\node[rel, anchor=north] at (\x,-0.22) {\v};}
\node[font=\scriptsize\itshape, text=gris50, anchor=north] at (11.35,-0.22) {(in progress)};
\end{tikzpicture}%
}
\caption[Milestone sequence of the project]{The six feature-complete milestones, from
the M1 walking skeleton to the M6 agentic console. Milestones M1--M4 each closed with a
semantically versioned release (v0.1.0--v0.4.0); M5 and M6 were in progress at the time
of writing.}
\label{fig:milestones}
\end{figure}

Three rules give the flow its discipline. First, work is \emph{issue-first}: no
feature is written before a board card exists for it, so every change is traceable
to a requirement. Second, every commit follows the Conventional Commits convention
\cite{conventionalcommits}, and every AI-assisted commit carries a
\texttt{Co-Authored-By: Claude} trailer, making AI involvement auditable from the
git history alone. Third, integration is governed by a \emph{CI-gated auto-merge}
(ADR-0003): a pull request merges only once the required checks are green, after
which the author reviews it asynchronously and retains full revert authority.
Figure~\ref{fig:pipeline} shows this pipeline.

\begin{figure}[htb]
\centering
\resizebox{\textwidth}{!}{%
\begin{tikzpicture}[
  >={Stealth[length=2.6mm]},
  stage/.style={draw=gris50, thick, rounded corners, align=center, text width=2.15cm, minimum height=1.25cm, fill=gris95, font=\scriptsize, inner sep=3pt},
  gate/.style={stage, draw=rojo, very thick, fill=rojo!8},
]
\node[stage] (issue) at (0,0)    {\bft{Issue}\\board card (FR-xx)};
\node[stage] (branch) at (2.7,0) {\bft{Branch}\\feature branch};
\node[stage] (pr) at (5.4,0)     {\bft{Pull request}\\Closes \#n};
\node[gate] (ci) at (8.1,0)      {\bft{CI gate}\\mypy, ruff, tests, coverage, xenon};
\node[stage] (merge) at (10.8,0) {\bft{Auto-merge}\\squash};
\node[stage] (val) at (13.5,0)   {\bft{Validate}\\author, async};
\foreach \a/\b in {issue/branch, branch/pr, pr/ci, ci/merge, merge/val} {\draw[->, gris50, thick] (\a) -- (\b);}
\end{tikzpicture}%
}
\caption[The continuous-integration development pipeline]{The development pipeline.
Each change starts from a board issue and reaches the main branch only through a
green CI gate (highlighted); the author's validation is asynchronous, after the
squash auto-merge, and may schedule follow-up work.}
\label{fig:pipeline}
\end{figure}

Two further practices keep an AI-driven codebase healthy over time. The Definition
of Done requires code, tests, docstrings and any affected diagrams to be updated
together; and a longitudinal quality audit is run before every milestone release to
catch the pathologies specific to AI-assisted development --- duplication, dead code,
complexity creep and pattern drift --- which per-change review cannot see.

\section{Technologies and tooling}
\label{sec:tech}

\subsection{Application stack}
\label{subsec:stack}

The system is a local, single-process web application. Table~\ref{tab:stack}
summarises its stack. The backend is written in Python~3.12, managed with the
\texttt{uv} toolchain, and exposes a FastAPI service that also serves the
single-page front end. The language-model layer is built on Pydantic~AI, which
provides the model-agnostic abstraction that lets the same code run against the
Claude API or a local Ollama model with only a configuration change. The retest
agent executes its commands inside an egress-locked Docker sandbox, and the front end
is a React single-page application built with Vite and Tailwind.

\begin{table}[htb]
\footnotesize
\renewcommand{\arraystretch}{1.3}
\begin{center}
\begin{tabular}{|L{2.8cm}|L{6.3cm}|L{3.4cm}|}
\hline
\rowcolor{tema!10}
\bft{Layer} & \bft{Technology} & \bft{Role} \\ \hline
Backend & Python 3.12 (\texttt{uv}), FastAPI, SQLAlchemy / SQLite & API, persistence, background jobs \\ \hline
Language model & Pydantic AI (Claude API; Ollama, local) & Finding extraction and the agentic retest \\ \hline
Report ingestion & pdfplumber & Parsing PDF reports into text \\ \hline
Sandbox & Docker (egress-locked network) & Isolated execution of gated commands \\ \hline
Front end & React 19, Vite, TypeScript, Tailwind CSS 4, React Query & Single-page application (FR-11) \\ \hline
Quality gates & ruff, mypy \texttt{-{}-strict}, pytest, coverage, xenon & Lint, types, tests, complexity \\ \hline
\end{tabular}
\end{center}
\caption[Technology stack of the system]{The technology stack, by layer. The
Pydantic AI abstraction is what makes the language-model backend swappable between a
cloud and a local provider without code changes.}
\label{tab:stack}
\end{table}

\subsection{Quality assurance}
\label{subsec:qa}

Quality is enforced mechanically rather than by convention. The continuous-integration
gate requires the strict type checker (mypy) and the linter/formatter (ruff) to pass,
test coverage to stay at or above 80\% of the source, and the cyclomatic-complexity
checker (xenon) to reject any function above its threshold --- a function that trips
the gate is refactored, never suppressed. Tests follow the classic pyramid: fast
unit tests with no input/output (language-model calls stubbed with Pydantic~AI's
test models), integration tests that wire real components together, and a small
number of dockerised system tests against the lab targets.

\subsection{The AI development environment}
\label{subsec:ai-env}

Beyond the application stack, the project relies on a bespoke development environment
built around Claude Code, the AI assistant. Repeated tasks are encoded as nine
project-specific \emph{skills} (for example, creating an ADR, adding a requirement,
opening a board-ready issue, or building the thesis), and quality checks as three
\emph{agents} (a whole-repository sanity auditor, a documentation curator, and a
thesis reviewer). Three \emph{hooks} automate governance directly in the editing
loop: one reminds the developer to open a board issue before coding, one formats
code on every edit, and one logs each assistant session to the public audit trail.
The assistant is further extended with a set of enabled plugins (among them GitHub,
Semgrep security scanning, Pydantic~AI support and up-to-date library documentation),
and the whole workflow is automated through seven GitHub Actions workflows: the
Kanban board automation, the CI gate, the documentation build, the pre-milestone
sanity audit, the security scan, the nightly system tests against the lab, and the
thesis build.

\subsection{Documentation as code}
\label{subsec:docs}

The project treats documentation as code: nothing that can be generated is
maintained by hand. The API reference is generated from the source docstrings with
mkdocstrings, and the UML class and package diagrams are generated from the code with
pyreverse, so neither can drift from the implementation. The architecture (a C4
model) and the sequence diagrams are authored in Mermaid inside the Markdown sources
and, by rule, updated in the same change as the code they describe. The whole set is
assembled into a static site with MkDocs --- built under a strict mode that fails on
any broken reference --- and published to GitHub Pages by a dedicated workflow, so
the public documentation always reflects the merged code. Documentation is a
first-class part of the Definition of Done: a change is complete only when its
docstrings, generated pages and affected diagrams are updated together, and the
\emph{doc-curator} agent audits for documentation drift as part of review.

\section{Use of Artificial Intelligence Tools}
\label{sec:ai-declaration}

In compliance with §6 of the ESII TFG regulation, this section declares the use of
artificial-intelligence tools in the production of this work. It is compiled from
the project's public audit trail: the curated usage log under
\texttt{docs/ai-usage/} --- maintained in real time through the early milestones and
completed retrospectively for the later ones from the automatically recorded session
files and the git history --- together with the \texttt{Co-Authored-By: Claude}
trailers in that history. Seventy-four of the eighty-two commits carry the trailer;
the remaining eight are author-only commits, and, because the project integrates by
squash-merge, each merge commit inherits the trailers of the assistant-authored work
it absorbs.

\paragraph{Tools used.} The single AI tool used for development was Claude Code,
Anthropic's agentic command-line assistant, driven primarily by the Fable~5 model
family during the early milestones and the Opus~4.x family thereafter, including for
assistance with this dissertation. Local open-weight models (Ollama, \texttt{qwen}
family) were used only \emph{as a component and comparison condition of the system
under study}, never to author the thesis or the code.

\paragraph{Type of use.} The assistance spanned six categories, in every case under
the author's direction and subject to the author's review:
\begin{itemize}
\item \emph{Research and analysis} --- reading the regulations and the proposal and
  surveying related work.
\item \emph{Design assistance} --- drafting the development plan and the
  thirty-five ADRs for the author's decision; the decisions themselves were the
  author's.
\item \emph{Code generation} --- implementing the backend (\texttt{src/revalid}) and
  the React front end (\texttt{frontend/src}) from the author's specifications.
\item \emph{Test writing} --- unit, integration and system tests across the pyramid.
\item \emph{Documentation generation} --- API documentation, the C4 architecture
  and sequence diagrams, and the project roadmap.
\item \emph{Thesis writing assistance} --- drafting and editing chapters (notably
  the State of the Art, the Introduction and this Methodology), including verifying
  that every cited reference resolves to a real source.
\end{itemize}

\paragraph{Affected areas.} In the repository, assistance touched the backend and
front-end source, the test suites, the ADRs and requirements, the documentation and
CI configuration, and the assistant's own skills, agents and hooks. In this
dissertation, the chapters listed above were drafted with assistance; all were
reviewed, corrected and approved by the author.

\paragraph{Authorship statement.} All design, architecture and scope decisions were
made by the author, as evidenced by the thirty-five ADRs and the requirements
specification that precede the code. Every AI output --- code, tests, documentation
and thesis prose --- was reviewed and validated by the author before acceptance,
through pull-request validation and the continuous-integration gate. The author is
the sole author of, and is responsible for, this work; the AI tool served as an
assistant that accelerated implementation, not as a substitute for the author's
judgement. The complete audit trail is public in the project repository under
\texttt{docs/ai-usage/}.
